% --- 1. GEOMETRÍA E IDIOMA ---
\usepackage[a4paper,left=2.75cm,right=2.6cm,top=2.5cm,bottom=2.5cm,heightrounded]{geometry}
\usepackage{emptypage}
\usepackage[english]{babel}
\usepackage{csquotes}
\usepackage{xspace}

% --- 2. COLORES Y GRÁFICOS (Antes de TikZ para evitar Option Clash) ---
\PassOptionsToPackage{dvipsnames}{xcolor}
\usepackage[dvipsnames]{xcolor}
\definecolor{goban}{HTML}{efce8f}
\definecolor{othello}{HTML}{485d3f}
\definecolor{linkdark}{HTML}{000055}

\usepackage{graphicx}
\usepackage{wrapfig}
\usepackage{float}
\usepackage{pdfpages}

% --- 3. MATEMÁTICAS Y TIPOGRAFÍA ---
\usepackage{mathtools}
\usepackage{amsmath,amsthm,amsfonts}
\usepackage{newpxtext,newpxmath}
\usepackage[no-math]{fontspec}
\usepackage{xeCJK}
\usepackage{microtype}

\linespread{1.03}
\setlength{\parskip}{0.5ex plus 0.2ex minus 0.2ex}
\setlength{\parindent}{1.2em}

\setCJKmainfont{Noto Serif CJK SC}
\setCJKmonofont{Noto Sans CJK SC}
\newfontfamily\cjkfont{Noto Serif CJK TC}[UprightFont=* Light, BoldFont=* Bold]

% --- 4. TIKZ Y PGFPLOTS ---
\usepackage{tikz}
\usepackage{tikz-cd}
\usepackage{pgfplots}
\pgfplotsset{compat=1.18}
\usetikzlibrary{babel, arrows.meta, bending, positioning, shadows, calc, shapes.geometric, decorations.text}

% --- 5. TABLAS Y DISEÑO EDITORIAL ---
\usepackage{booktabs}
\usepackage{multicol}
\usepackage{multirow}
\usepackage{epigraph}
\usepackage[margin=1cm, font=small, labelfont=bf, labelsep=endash]{caption}

% --- 6. BIBLIOGRAFÍA (BibLaTeX + Biber) ---
\usepackage[
backend=biber,
style=numeric-comp,
natbib=true,
url=false,
doi=true,
eprint=true,
isbn=false,
maxcitenames=2,
maxbibnames=99,
giveninits=true,
sorting=nty
]{biblatex}
\usepackage{xurl}

\renewcommand*{\bibfont}{\small}
\addbibresource{bibliography/biblio.bib}

\defbibheading{lineabib}[\bibname]{%
	\chapter*{#1}%
	\markboth{#1}{#1}%
	\nopagebreak\vspace{-4ex}\noindent\rule{\linewidth}{1.2pt}\par
}

% --- 7. ÍNDICE GENERAL (ToC) ---
\usepackage{tocloft}
\setcounter{tocdepth}{3}
\renewcommand{\cftchapfont}{\Large\bfseries}
\renewcommand{\cftaftertoctitle}{\par\noindent\rule{\linewidth}{1.2pt}}
\renewcommand{\cftchappagefont}{\Large\bfseries}
\renewcommand{\cftchapleader}{\hfill}
\renewcommand{\cftchappresnum}{Part }
\setlength{\cftchapnumwidth}{6.25em}
\setlength{\cftsecnumwidth}{2.3em}
\setlength{\cftsubsecnumwidth}{3.5em}
\setlength{\cftbeforetoctitleskip}{-0.5cm}

% --- 8. FORMATO DE CAPÍTULOS ---
\renewcommand{\thechapter}{\Roman{chapter}}
\renewcommand{\thesection}{\arabic{chapter}.\arabic{section}}
\renewcommand{\thefigure}{\arabic{chapter}.\arabic{figure}}
\renewcommand{\thetable}{\arabic{chapter}.\arabic{table}}

\makeatletter
\renewcommand{\@makechapterhead}[1]{%
	{\parindent \z@
		\normalfont
		\Huge\bfseries \Roman{chapter}: \nobreak
		\Huge #1\par\nobreak
		\vskip 20pt
	}%
}
\makeatother

% --- 9. ENCABEZADOS Y PIES DE PÁGINA ---
\usepackage{fancyhdr}
\setlength{\headheight}{25pt}
\pagestyle{fancy}
\fancyhf{}
\fancyhead[R]{\bfseries\thepage}
\fancyhead[L]{\small\scshape\nouppercase{\leftmark}}
\renewcommand{\headrulewidth}{0.5pt}
\renewcommand{\footrulewidth}{0pt}
\renewcommand{\chaptermark}[1]{%
	\markboth{Part \thechapter: #1}{}%
}

% --- 10. ENTORNOS MATEMÁTICOS ---
\newtheorem{theorem}{Theorem}[section]
\newtheorem{lemma}[theorem]{Lemma}
\newtheorem{corollary}[theorem]{Corollary}

\theoremstyle{definition}
\newtheorem{definition}{Definition}[section]

% Símbolo Q.E.D. 
\renewcommand{\qedsymbol}{\raisebox{-1.8pt}{\huge\(\bullet\)}}

% --- 11. HIPERVÍNCULOS Y METADATOS (Siempre al final) ---
\usepackage{hyperref}
\hypersetup{
	colorlinks=true,
	linkcolor=linkdark,
	citecolor=linkdark,
	urlcolor=linkdark,
	bookmarksnumbered=true,
	pdfpagemode=UseOutlines,
	pdfstartview=FitH,
	pdftitle={Wéi Dào - The Art of Surrounding},
	pdfauthor={W. Adrián Arias}
}